\section*{\centering ВСТУП}\addcontentsline{toc}{section}{ВСТУП}

\textbf{Обґрунтування вибору теми дослідження.} Цифрова трансформація освіти, що відбувається в умовах глобального інформаційного суспільства, формує нові вимоги до професійної діяльності педагога. Імерсивні технології (VR, AR, MR, XR) стають невід'ємним складником сучасного освітнього середовища, забезпечуючи занурення здобувачів освіти у навчальний зміст і розширюючи можливості педагогічної діяльності. Європейська Комісія у стратегічних документах рекомендує підтримувати дослідження, розробку та перевірку цифрових освітніх засобів і новітніх технологій, у тому числі розширеної, віртуальної та доповненої реальності, а також запровадити обов'язкову складову професійної підготовки майбутніх учителів~-- <<ІКТ в освіті>>.

За таких умов учитель інформатики має бути не лише користувачем готових цифрових рішень, а й творцем інноваційних освітніх середовищ, здатним інтегрувати технології доповненої, віртуальної та змішаної реальності у власну педагогічну практику. Проте аналіз стану підготовки майбутніх учителів інформатики в Україні свідчить про відсутність усталених методик формування готовності до створення імерсивних освітніх ресурсів (ІОР), невизначеність критеріїв їхньої педагогічної ефективності та недостатню інтеграцію психолого-педагогічних підходів до цифрового дизайну.

Актуальність дослідження зумовлена суперечностями між: зростаючою потребою освітньої практики в імерсивних ресурсах та недостатнім рівнем готовності майбутніх учителів інформатики до їх розробки; наявним технологічним потенціалом WebAR/VR/XR та відсутністю науково обґрунтованої методики навчання розробки ІОР; необхідністю інтеграції імерсивних технологій в освітній процес та нерозробленістю організаційно-педагогічних умов такої інтеграції у систему підготовки вчителів інформатики.

\textbf{Зв'язок роботи з науковими програмами, планами, темами.} Дисертацію виконано у межах фундаментального дослідження <<Теоретико-методичні засади проєктування імерсивного хмаро орієнтованого освітнього середовища університету>> (державний реєстраційний номер 0121U113711). Тему дисертації затверджено вченою радою Криворізького державного педагогічного університету.

\textbf{Об'єкт дослідження}~-- процес професійної підготовки майбутніх учителів інформатики у закладах вищої освіти.

\textbf{Предмет дослідження}~-- методика навчання розробки імерсивних освітніх ресурсів майбутніх учителів інформатики.

\textbf{Мета дослідження}~-- теоретично обґрунтувати та експериментально перевірити методику навчання розробки імерсивних освітніх ресурсів майбутніх учителів інформатики.

\textbf{Гіпотеза дослідження.} Ефективність методики навчання розробки імерсивних освітніх ресурсів майбутніх учителів інформатики підвищиться, якщо: визначити теоретичні засади навчання розробки ІОР; розробити та впровадити структурно-функціональну модель підготовки; забезпечити реалізацію педагогічних умов навчання розробки ІОР із використанням спеціально розробленого навчально-методичного забезпечення (навчальний посібник з WebAR, курс у LMS Moodle, лабораторний практикум).

\textbf{Завдання дослідження:}
\begin{enumerate}
    \item Визначити теоретичні засади навчання розробки імерсивних освітніх ресурсів у підготовці майбутніх учителів інформатики.
    \item Обґрунтувати організаційно-педагогічні основи навчання розробки ІОР: дидактичні підходи та принципи, структурно-функціональну модель підготовки, педагогічні умови.
    \item Організувати та провести дослідно-експериментальну роботу з перевірки ефективності розробленої методики навчання.
\end{enumerate}

\textbf{Методи дослідження.} Для розв'язання поставлених завдань використано комплекс методів у межах змішаного дослідницького дизайну (конвергентний паралельний дизайн):
\begin{itemize}
    \item \textit{теоретичні}: аналіз, синтез, систематизація та узагальнення наукових джерел з проблеми дослідження; порівняльний аналіз освітньо-професійних програм підготовки вчителів інформатики; моделювання структурно-функціональної моделі підготовки;
    \item \textit{емпіричні}: педагогічний експеримент (констатувальний, формувальний, контрольний етапи), анкетування, тестування, спостереження, експертне оцінювання; спрямований якісний контент-аналіз транскриптів відеолекцій;
    \item \textit{статистичні}: критерій $\chi^2$ Пірсона для перевірки статистичної значущості результатів експерименту;
    \item \textit{інтерпретаційні}: тріангуляція кількісних та якісних даних для верифікації педагогічних механізмів.
\end{itemize}

\textbf{Експериментальна база дослідження.} Дослідно-експериментальна робота здійснювалася на базі Криворізького державного педагогічного університету (кафедра інформатики та прикладної математики) протягом 2020--2025~рр. У дослідженні брали участь студенти спеціальності 014 Середня освіта (Інформатика).

\textbf{Наукова новизна одержаних результатів:}

\textit{уперше:}
\begin{itemize}
    \item визначено ІОР як педагогічний об'єкт у системі підготовки вчителів інформатики, уточнено сутнісні характеристики імерсивних освітніх ресурсів як інтерактивних засобів нового покоління, побудованих на моделюванні віртуальних або доповнених середовищ;
    \item обґрунтовано структурно-функціональну модель підготовки майбутніх учителів інформатики до розробки ІОР, що включає п'ять взаємопов'язаних блоків: мотиваційно-цільовий, теоретико-методологічний, змістовий, організаційно-технологічний, рефлексивно-оцінювальний;
    \item розроблено та апробовано трикрокову методику навчання WebAR з інтеграцією моделей машинного навчання (TensorFlow.js, Teachable Machine);
    \item виявлено темпоральну динаміку факторів когнітивно-афективної моделі імерсивного навчання (CAMIL) у процесі навчання розробки ІОР та верифіковано ефективність методики за допомогою конвергентного паралельного дизайну змішаних методів дослідження;
\end{itemize}

\textit{удосконалено:}
\begin{itemize}
    \item класифікацію імерсивних освітніх ресурсів за типом занурення, технологічною основою та дидактичним призначенням;
    \item методику навчання розробки WebAR-застосунків майбутніх учителів інформатики з використанням MindAR, Three.js та A-Frame;
    \item підходи до валідації методик навчання імерсивних технологій шляхом тріангуляції кількісних та якісних даних;
\end{itemize}

\textit{набуло подальшого розвитку:}
\begin{itemize}
    \item наукові положення щодо використання технологій доповненої та віртуальної реальності в освіті;
    \item підходи до інтеграції машинного навчання у WebAR-застосунки для освітніх цілей;
    \item підходи до якісного аналізу педагогічного дискурсу в навчанні імерсивних технологій (спрямований контент-аналіз лекційних транскриптів).
\end{itemize}

\textbf{Практичне значення одержаних результатів} полягає в розробці:
\begin{itemize}
    \item навчального посібника з WebAR (10 розділів: вступ до WebAR, інструменти розробки, принципи роботи AR у браузерах, відстеження зображень MindAR, відстеження обличчя, відстеження простору WebXR, інтеграція машинного навчання, інші технології WebAR, лабораторний практикум);
    \item курсу <<Синтетична реальність в освіті>> у LMS Moodle на платформі moodle.kdpu.edu.ua (18 тижневих модулів);
    \item лабораторного практикуму (10 варіантів творчих завдань із 3D-моделювання та AR);
    \item віртуальної хімічної лабораторії з доповненою реальністю (AR-застосунок для якісного аналізу).
\end{itemize}

\textbf{Особистий внесок здобувача.} У працях, опублікованих у співавторстві, здобувачу належать: порівняльний аналіз інструментів розробки WebAR (A-Frame, AR.js, Three.js, JSARToolKit, Babylon.js); розробка та апробація WebAR-квестів для профорієнтації на базі КДПУ; програмна реалізація WebAR-застосунків з інтеграцією MindAR та Three.js; розробка трикрокової методики навчання WebAR з моделями машинного навчання; створення та тестування WebAR-застосунків з TensorFlow.js handpose та Teachable Machine; програмна реалізація віртуальної хімічної лабораторії.

\textbf{Апробація результатів дослідження.} Основні положення та результати дослідження доповідалися та обговорювалися на наукових конференціях різного рівня:
\begin{itemize}
    \item \textit{Міжнародних}:
3rd International Workshop on Augmented Reality in Education (AREdu 2020) (Кривий Ріг, 2020),
XII International Conference on Mathematics, Science and Technology Education (ICon-MaSTEd 2020) (Кривий Ріг, 2020),
International Workshop on Computer Science \& Software Engineering at Schools and in the Workplace (CS\&SE@SW 2020) (Кривий Ріг, 2020),
9th Workshop on Cloud Technologies in Education (CTE 2021) (Кривий Ріг, 2021),
Digital Humanities Workshop (DHW 2021) (Київ, 2021),
11th Illia O. Teplytskyi Workshop on Computer Simulation in Education (CoSinE 2024) (Кривий Ріг, 2024).
\end{itemize}

\textbf{Публікації.} Основні результати дисертації відображено у 8~наукових працях (усі у співавторстві), у тому числі: 4~статті у наукових періодичних виданнях інших держав, включених до міжнародної наукометричної бази Scopus; 3~статті у наукових фахових виданнях України; 4~публікації засвідчують апробацію матеріалів дисертації; 1~публікація додатково відображає наукові результати дисертації.

\textbf{Структура та обсяг дисертації.} Дисертація складається з анотації, вступу, трьох розділів, висновків, списку використаних джерел (\arabic{citenum@totc}~найменувань) та 5~додатків. Загальний обсяг роботи~-- \pageref{last_page}~сторінок, з них основного тексту~-- \pageref{end_main_text}~сторінок. Робота містить \totalfigures~рисунків та \totaltables~таблиць.
